% --- Frontpage ----------------------------------------------------------------
\maketitle

% --- Notation & Prerequisites -------------------------------------------------
\chapter*{Notation}
\addcontentsline{toc}{chapter}{Notation}
\begin{itemize}
    \item $\ExtC$: Extended complex plane.
    \item $\Delta$: Unit disk in $\C$.
    \item $N$: Neighborhood of a point.
    \item $fg$: Composition of $f$ and $g$ ($f \circ g$).
    \item $f^n$: The composition of $f$ with itself $n$ times.
    \item $\Id = f^0$: Identity function.
    \item $f^{(n)}$: $n$-th derivative of $f$.
    \item $\zeta$: Fixed point.
    \item $R(z)$: Rational function in complex variable $z$.
    \item $\sigma$: Chordal metric.
    \item $\sigma_0$: Spheric metric.
    \item $\gamma$: A curve.
    \item $E(\gamma)$: Exterior component of a closed curve.
    \item $\Omega$: An interior component of a closed curve.
    \item $v_f(z_0)$: Valency of $f$ at $z_0$.
    \item $m(R, z_0)$: Multiplier of $R$ at $z_0$.
\end{itemize}

\chapter*{Prerequisites}
\addcontentsline{toc}{chapter}{Prerequisites}

\section*{From Complex Analysis}
\begin{itemize}
    \item $d$-fold map: We say $f$ is a $d$-fold map of $V$ onto $W$ if, for
    every $w \in W$, the equation $f(z) = w$ has exactly $d$ solutions in $V$
    counting multiple solutions (by their multiplicity).
    \item Maximum Modulus Theorem.
    \item Schwarz's Lemma.
    \item Rouché's Theorem.
    \item Hurwitz's Theorem: If a sequence of univalent analytic maps $f_n$ on a
    domain $D$ converges uniformly to $f$, then $f$ is either constant or
    univalent in $D$.
\end{itemize}

\section*{From Topology}
\begin{itemize}
    \item Metric spaces: Uniform continuity.
    \item Compactness.
    \item Connectedness.
\end{itemize}

% --- Table of Contents --------------------------------------------------------
\tableofcontents